\section{State of the Art}
% Organised around the problem: each subsection groups the studies by approach, then states agreements, trade-offs and implications.

% table* spans BOTH columns. Do not remove the asterisk.
\begin{table*}[!t]
\caption{Comparison of the reviewed literature}
\label{tab:literature}
\centering
\footnotesize
\setlength{\tabcolsep}{4pt}
\renewcommand{\arraystretch}{1.18}
\begin{tabular}{|L{0.11\textwidth}|L{0.11\textwidth}|L{0.18\textwidth}|L{0.18\textwidth}|L{0.17\textwidth}|L{0.15\textwidth}|}
\hline
\textbf{Study} & \textbf{Approach} & \textbf{Context and technology} & \textbf{Evidence / results} & \textbf{Limitations} & \textbf{Relevance to TrackSeed} \\ \hline
Chalhoub and Ayer~\cite{chalhoub2019} & Head-mounted AR layout & Indoor electrical point layout; HoloLens with paper marker & About 5~cm error; smaller near the marker & Indoor, small area; no time comparison & Accuracy depends on distance from the reference \\ \hline
Kim and Olsen~\cite{kim2025isarc}, \cite{kim2024ahfe} & Handheld, head-mounted and GNSS-assisted AR & Outdoor utility placement; 32 students; total-station ground truth & Mean error 8.3~cm (HoloLens~2) to 15.9~cm (iPhone) & Small area; 8 observations per device and pipe; students & Phone accuracy is close to the FS tolerance \\ \hline
Feigl et al.~\cite{feigl2020} & Visual-inertial tracking & ARCore, ARKit, HoloLens in a 1{,}600~m\textsuperscript{2} hall; optical ground truth & About 17~m error per 120~m; scale error up to 14.4~cm/m & Dynamic indoor industrial setting & Re-alignment is needed over long distances \\ \hline
Scargill et al.~\cite{scargill2021} & Visual-inertial tracking & Three phones, more than 200 indoor scenes & Drift of a few cm; 4.8~cm (ARCore) vs 25.1~cm (ARKit) after walking away & Room scale; preprint & Low texture and light increase drift (asphalt risk) \\ \hline
Nunes et al.~\cite{nunes2023} & Multi-marker registration & Large indoor industrial space; HoloLens with several markers & Two markers significantly reduce angular error; more points improve accuracy & Indoor; head-mounted device & Supports two-point alignment at the start \\ \hline
Jiang et al.~\cite{jiang2023} & Smartphone RTK GNSS & Unmodified phones; static and kinematic tests in open sky & 1.1--2.1~cm static, about 4~cm kinematic RMS error & Algorithm study; no AR; open sky only & Absolute anchoring without drift \\ \hline
Ajayi et al.~\cite{ajayi2023} & Pub/sub vs client--server & MQTT vs WebSocket for IoT telemetry & MQTT: higher throughput, fewer resources; WebSocket: lower latency & Single controlled test; telemetry data & Protocol choice for phone--dashboard sync \\ \hline
Chen and Ran~\cite{chen2019}, Ignatov et al.~\cite{ignatov2019} & Edge vs cloud inference & Review of edge deep learning; smartphone AI benchmark & On-device avoids network latency and privacy risks; mobile NPUs approach desktop GPUs & Review and benchmark, not this task & Run the cone check on the phone \\ \hline
V\"odisch et al.~\cite{vodisch2022}, Dhall et al.~\cite{dhall2019} & FS cone detection & Shared cone dataset; real-time detection on a driverless car & Annotated FS cone data; real-time detection on embedded hardware & Car-mounted cameras, not handheld & Training data and model baseline \\ \hline
Yang et al.~\cite{yang2015}, Yosinski et al.~\cite{yosinski2014}, Sandler et al.~\cite{sandler2018} & Fine-grained car classification and transfer learning & Large car dataset; transfer of pretrained features; mobile CNN & Accuracy depends on viewpoint; pretrained features help small datasets; 72\% ImageNet top-1 with 3.4M parameters & Public car models, not a team's own cars & Self-made dataset with pretrained mobile backbone \\ \hline
\end{tabular}
\end{table*}

\subsection{AR for on-site layout and spatial guidance}
Construction layout, where workers mark the positions of building elements from a design model, is the task closest to setting out a track, and it is the domain where AR placement has been measured most directly. The studies found differ mainly in the device used and in the size and type of the test area.

\textit{Head-mounted AR in small indoor areas.} Chalhoub and Ayer asked industry practitioners to lay out electrical points with a HoloLens registered to a paper marker, and measured each point with a laser distance meter~\cite{chalhoub2019}. Content was displayed about 5~cm from its design position, and the error was smaller near the marker. Showing more content, or content of different types, also reduced accuracy in elevation. The test was carried out indoors, over an area much smaller than a track.

\textit{Head-mounted, handheld and GNSS-assisted AR outdoors.} Kim and Olsen compared four devices in an outdoor, partly shaded simulation of underground utility placement, with 32 construction students placing survey prisms where they perceived each virtual pipe and a robotic total station recording the positions~\cite{kim2025isarc}. The mean error was 3.28~in (about 8.3~cm) with a HoloLens~2, 3.52~in (8.9~cm) with Trimble SiteVision, which adds a GNSS receiver, 3.59~in (9.1~cm) with an iPad Pro and 6.26~in (15.9~cm) with an iPhone~13 Pro Max, with eight observations per device and pipe. An earlier report on the same scenario also found the HoloLens to be the most accurate device and concluded that the devices did not meet the strict tolerances of construction work~\cite{kim2024ahfe}.

The two groups of studies agree that AR placement reaches centimetre-to-decimetre accuracy over small areas, that head-mounted devices are more accurate than handheld ones, and that accuracy depends on the distance from the reference used to register the content. They differ by a factor of about two in the error they report, which can be explained by the indoor versus outdoor setting and by the participants' experience, so the values cannot be compared directly. For TrackSeed, the most relevant figure is the phone result: an error of about 16~cm outdoors is close to the provisional tolerance of 10--15~cm assumed for an FS track, and it was obtained over a much smaller area than a track. The literature therefore supports AR guidance as feasible, but not as accurate enough by default; the accuracy of phone-based setout over a full track is the main hypothesis TrackSeed has to test.

\subsection{Mobile AR tracking accuracy and alignment to the ground}
Placement accuracy over a full track depends on how well the phone keeps its position as the user walks, and on how the virtual plan is aligned to the real ground. Four approaches appear in the literature: relying on the phone's own visual-inertial tracking after an initial alignment, registering the content repeatedly against several references, localising the phone against a visual map, and using satellite positioning.

\textit{Visual-inertial tracking alone.} ARCore estimates the phone's pose from camera features combined with inertial measurements~\cite{arcore}. In a 1{,}600~m\textsuperscript{2} industrial hall with walking distances of 60~m, Feigl et al.\ found that ARCore, ARKit and HoloLens accumulated an average error of about 17~m per 120~m walked, with a scale error of up to 14.4~cm/m that grew roughly in proportion to the path length~\cite{feigl2020}. At room scale, Scargill et al.\ measured the drift of virtual objects on three phones in more than 200 environments and found errors of a few centimetres that depended on what the user did between views: when the user walked away and returned, the mean drift was 4.8~cm with ARCore and 25.1~cm with ARKit~\cite{scargill2021}. Drift also increased in dim scenes and in scenes with little texture or low contrast.

\textit{Registration against several references.} Nunes et al.\ proposed a calibration model that registers AR content in a large indoor industrial space with several markers, tested on a HoloLens; they report that single small markers are numerically unstable, that using two markers significantly reduces the angular discrepancy between real and virtual points, and that accuracy improves as more calibration points are added~\cite{nunes2023}.

\textit{Visual positioning.} The ARCore Geospatial API localises the phone against imagery from Google Street View and is described by Google as more accurate than GPS alone, but the documentation does not give accuracy figures and the service only works where Street View coverage exists~\cite{arcoregeo}. Skabek et al.\ compared it with other AR localisation platforms for indoor and outdoor campus navigation~\cite{skabek2026}, a use case whose accuracy needs are far looser than cone placement. No independent evaluation at the centimetre level was found.

\textit{Satellite positioning.} Standard phone GNSS is not accurate enough for cone placement, but RTK corrections change this: Jiang et al.\ reached root-mean-square errors of 1.1--2.1~cm in a static test and about 4~cm in a kinematic test with unmodified smartphones in open-sky conditions~\cite{jiang2023}. In the outdoor study of Kim and Olsen, the GNSS-assisted SiteVision was more accurate than both handheld devices but not better than the HoloLens~\cite{kim2025isarc}.

The two tracking studies appear to disagree, with errors of metres in one and centimetres in the other, but the difference is explained by scale: the error of visual-inertial tracking grows with the distance walked~\cite{feigl2020}, and a room-scale study cannot reveal it. Both agree that the environment matters, and Scargill et al.'s finding that low texture increases drift is a risk for the uniform asphalt of a test site, although no study measured it on tarmac. Registration against several references addresses the angular error of a single marker~\cite{nunes2023}, but it has only been tested indoors with a head-mounted device. The trade-off is clear: visual tracking is free and works anywhere but drifts over distance; visual positioning depends on map coverage and lacks published accuracy; RTK GNSS gives centimetre accuracy without drift but needs a receiver, a correction source and an open sky, which an FS test site usually provides. For TrackSeed, these findings mean that a single alignment at the start line cannot keep a track of several hundred metres within tolerance, and that the plan must be re-aligned at reference points along the track or anchored to absolute positions.

\subsection{Real-time platforms and synchronisation}
Live updates between the phone, the cloud and the dashboard can follow a client--server model, typically WebSocket connections, or a publish--subscribe model such as MQTT~\cite{ajayi2023}.\ A comparison between both for streaming IoT telemetry and found that MQTT achieved higher throughput and used fewer system resources, while WebSocket had slightly lower round-trip times~\cite{ajayi2023}. The evidence found for this theme is limited to a single controlled comparison with telemetry data, so it shows a trade-off rather than a general ranking. For TrackSeed, the volume of data is small (one message per confirmed cone) and the receiving end is a web dashboard, which favours WebSocket's direct browser support and lower latency over MQTT's efficiency at scale; this is assessed further in Section~III.

\subsection{Edge AI for cone verification and car identification}
TrackSeed runs two vision models, one that checks each placed cone and one that identifies the team car being used; each which can run in the cloud, on the phone or in a hybrid of both.The trade-off being: sending camera frames to the cloud gives access to more computing power but adds latency and privacy risks while using uplink bandwidth, while on-device inference avoids these costs but is limited by the device's computing power and memory~\cite{chen2019}. Frameworks such as TensorFlow Lite (now LiteRT) provide on-device inference and mobile hardware has narrowed the gap: in a benchmark of smartphone AI accelerators, newest mobile neural processing units were approaching the performance of CUDA-compatible Nvidia graphics cards~\cite{ignatov2019}.

No public dataset contains a team's own cars, so TrackSeed must use a small dataset made by the authors. TrackSeed has only a few classes, but cars from the same team often look alike. Its images also all come from one site, so training and testing on frames from the same day would overestimate accuracy. TrackSeed therefore fine-tunes a pretrained mobile network on photos of each car from all viewpoints, taken on several test days.

\subsection{Comparative synthesis}
Table~\ref{tab:literature} summarises the studies reviewed. Three findings stand out. First, every study that measured position over distance found that error grows with the distance from the reference, whether that reference is a marker~\cite{chalhoub2019} or the starting point of tracking~\cite{feigl2020}; the studies disagree on magnitude only because they test different scales. Second, almost all accuracy evidence comes from areas of a few metres to tens of metres, while an FS track covers hundreds; the only large-scale study measured errors far above any useful tolerance without re-alignment~\cite{feigl2020}. Third, phones are consistently less accurate than head-mounted devices~\cite{kim2025isarc}, yet phones are the only devices FS teams have available. Fourth, pretrained mobile networks make on-device car classification feasible with little data~\cite{yosinski2014}, \cite{sandler2018}, but the reported accuracies come from large public datasets, not from a few similar cars photographed at one site. What cannot be concluded from the literature is whether phone-based AR, combined with re-alignment or RTK positioning, keeps a full track within tolerance: no study has tested this. No work was found that combines placement guidance, automatic verification of each placed object, identification of the vehicle under test and a live two-way digital twin, which is the gap TrackSeed addresses.
